% =============================================================================
% Standalone chapter / appendix: TAC simulations
% -----------------------------------------------------------------------------
% Intended use:
%   - As an appendix (or chapter) in Resilience Report V8.0
%   - Alone: compile tac_simulations_standalone.tex
%
% Figures are copies of knitr PNGs from report/html/*_files/figure-html/,
% staged under report/latex/figs/. Refresh them (all in R) with:
%   source("report/finalise/02_stage_figures.R"); stageFigures()
% =============================================================================

\chapter{Simulation of future TACs}
\label{ch:tac-sim}

\paragraph{Plain-language framing.}
This appendix explores \emph{what catch and biomass paths might look like}
if recent fishing pressure (or alternative $F_{\mathrm{MSY}}$ targets) were
held forward under different recruitment or productivity assumptions. It is
\textbf{not} an ICES or ICCAT advice product, and it does not recommend TACs
for the coming year. The aim is to support the resilience narrative in the
main report: how sensitive near-term catches and stock status are when
productivity is poorer than the recent mean, or when the latest recruitment
regime persists.

Technically, the appendix summarises deterministic short- to medium-term Total
Allowable Catch (TAC) simulations for the stocks retained in the BIM
resilience analysis. Calculations live in the companion R Markdown workflows
(\texttt{01\_pelagics.Rmd}, \texttt{02\_demersal.Rmd}, \texttt{03\_nephrops.Rmd},
\texttt{04\_iccat.Rmd}). Here we document the shared scenario design, show the
principal diagnostic and forecast figures, and point to the tabulated catch
paths under \texttt{data/TAC/csv/}. Broader climate synthesis and
cross-fishery conclusions belong in the parent report text.

\paragraph{Key messages for non-modelling readers.}
\begin{itemize}
  \item Compare scenarios \emph{within} a stock first: the vertical gap between
        lines is the effect of the assumption that changed.
  \item Status-quo $F$ with lower recruitment asks: ``If we keep fishing as we
        do now, and productivity worsens, what happens to catch and biomass?''
  \item $F_{\mathrm{MSY}}$ paths ask the same productivity questions under a
        different fishing intensity.
  \item Regime paths ask: ``What if the current recruitment regime continues?''
  \item Use $B/B_{\mathrm{trigger}}$ (in the CSVs) to judge status; absolute
        catch scales differ widely between stocks.
\end{itemize}

% -----------------------------------------------------------------------------
\section{Shared design}
\label{sec:tac-design}

\subsection{Workflow}
\label{sec:tac-workflow}

For each stock group the analysis follows three steps:

\begin{enumerate}
  \item \textbf{Assemble operating models (OMs)} from ICES Stock Assessment
        Graphs (SAG) series, accepted assessment objects (SS3, SAM,
        \texttt{icesdata}), and---for Nephrops---JABBA surplus-production fits.
  \item \textbf{Bridge to the advice year} with agreed catch advice, so that
        the projections start from a common near-term state rather than the
        raw terminal assessment year.
  \item \textbf{Project forward} under a fixed control table of fishing
        mortality (or harvest) targets and recruitment / process-error
        multipliers, then export catch and biomass trajectories to CSV.
\end{enumerate}

ICES age-structured stocks use \texttt{FLasher::fwd} with stock--recruit
relationships fitted via Beverton--Holt equilibria (\texttt{calcEq}). Nephrops
functional units use biodynamic OMs conditioned on JABBA, projected with
\texttt{mpb::fwd}. North Atlantic albacore uses the ICCAT SS3 OM with the same
status-quo $F$ target under scaled recruitment deviations.

\subsection{Scenarios}
\label{sec:tac-scenarios}

Age-structured ICES stocks (pelagic and demersal) share one forecast control
table. For each row, Table~\ref{tab:tac-scen} states what is held constant,
what is changed, and the question the scenario is meant to answer.
Status-quo $F$ is the mean of $\bar{F}$ over the last three assessment years.
Age-structured projections (pelagic, demersal, albacore) run through 2050;
the Nephrops harvest-rate projections run through 2041.

\begin{table}[htbp]
\centering
\caption{Age-structured ICES TAC scenarios: constant controls, changed
inputs, and intended question.}
\label{tab:tac-scen}
\small
\begin{tabular}{@{}p{0.18\textwidth}p{0.24\textwidth}p{0.24\textwidth}p{0.28\textwidth}@{}}
\toprule
Scenario & Held constant & Changed & Question answered \\
\midrule
FStatus Quo
  & Recent mean $F$; recent recruitment residual mean
  & --- (reference path)
  & What is the baseline catch/biomass path if recent fishing and recent
    productivity continue? \\
\addlinespace
FStatus Quo 75\%
  & Recent mean $F$
  & Recruitment residual $\times 0.75$
  & If fishing stays as now, how much do catch and status erode when
    recruitment is moderately poorer? \\
\addlinespace
FStatus Quo 50\%
  & Recent mean $F$
  & Recruitment residual $\times 0.50$
  & Same as above under a stronger productivity shock. \\
\addlinespace
FMSY
  & $F_{\mathrm{MSY}}$; recent residual mean
  & Fishing intensity set to $F_{\mathrm{MSY}}$ (vs status quo)
  & Relative to status quo, does fishing at $F_{\mathrm{MSY}}$ improve or
    worsen biomass under recent productivity? \\
\addlinespace
Regime
  & $F_{\mathrm{MSY}}$
  & Latest stock--recruit residual \emph{regime} replaces the short-window
    residual mean
  & If the current recruitment regime persists, how do $F_{\mathrm{MSY}}$
    catch and biomass paths compare with the recent-mean residual case? \\
\bottomrule
\end{tabular}
\end{table}

Nephrops scenarios are formulated in harvest-rate space and follow the same
logic, but the productivity shock is applied through the biodyn process-error
multiplier rather than a recruitment residual. The scenario labels in
\texttt{neph-f.csv} are \texttt{FStatus Quo} (status-quo harvest rate, process
error $= 1$; the reference path), \texttt{FStatus Quo 75\%} and
\texttt{FStatus Quo 50\%} (same status-quo harvest, process-error multiplier
$0.75$ and $0.50$; i.e.\ fishing unchanged while productivity is progressively
poorer), \texttt{Fmsy} ($F_{\mathrm{MSY}}$ harvest), and \texttt{PE}
($F_{\mathrm{MSY}}$ with autocorrelated process-error noise). Note the
capitalisation differs from the age-structured runs (\texttt{Fmsy} here versus
\texttt{FMSY} for pelagic/demersal), and the Nephrops set has no \texttt{Regime}
row. Albacore reports status-quo $F$ under recruitment multipliers 1, 0.75 and
0.5 (reference, moderate, and strong productivity reduction).

\subsection{Exports}
\label{sec:tac-exports}

Harmonised forecast tables are written as
\texttt{data/TAC/csv/\{pel,dem,neph,iccat\}-f.csv} with columns \texttt{Scenario},
\texttt{sid}, \texttt{Year}, \texttt{Catch}, \texttt{B}, and \texttt{Btrig}
($B/B_{\mathrm{trigger}}$). Full trajectory objects are stored beside them as
\texttt{*-fct.RData}.

% -----------------------------------------------------------------------------
\section{Pelagic stocks}
\label{sec:tac-pel}

The pelagic group comprises boarfish (\texttt{boc.27.6-8}), western horse
mackerel (\texttt{hom.27.2a3a4a5b6a7a-ce-k8}), Northeast Atlantic mackerel
(\texttt{mac.27.nea}), and blue whiting (\texttt{whb.27.1-91214}). Assessments
are taken from SS3, the 2025 mackerel FLStock, and SAM, then reconciled to
ICES SAG time series before the advice bridge
(Figure~\ref{fig:pel-sag}).
For resilience reading: compare status-quo paths under reduced recruitment
with the Regime path to see whether recent productivity already implies
materially lower long-run catch than the multi-year residual mean.

\begin{figure}[htbp]
  \centering
  \includegraphics[width=0.95\textwidth]{pel-sag-vs-stock.png}
  \caption{Pelagic stocks: SAG summaries versus FLStock metrics scaled
           relative to MSY reference points where available.}
  \label{fig:pel-sag}
\end{figure}

Recruitment residual diagnostics and regime polygons underpin the Regime
scenario (Figure~\ref{fig:pel-rec}). Projected catch under the five control
rows is shown in Figure~\ref{fig:pel-fcst}; tabulated paths are in
\texttt{pel-f.csv}.

\begin{figure}[htbp]
  \centering
  \includegraphics[width=0.95\textwidth]{pel-rec-residuals.png}
  \caption{Pelagic stock--recruit residuals. Shaded polygons mark regime
           segments used to define the Regime recruitment multiplier.}
  \label{fig:pel-rec}
\end{figure}

\begin{figure}[htbp]
  \centering
  \includegraphics[width=0.95\textwidth]{pel-forecast-catch.png}
  \caption{Pelagic forecast catch by stock and scenario (status quo $F$,
           reduced recruitment under status quo $F$, $F_{\mathrm{MSY}}$, and
           Regime).}
  \label{fig:pel-fcst}
\end{figure}

% -----------------------------------------------------------------------------
\section{Demersal stocks}
\label{sec:tac-dem}

The demersal suite covers anglerfishes, northern hake, megrims, and whitings
in the west and northern shelf assessment areas. Inputs combine
\texttt{icesdata}, SAM, and SS3 (with season aggregation where required).
Alignment against SAG is shown in Figure~\ref{fig:dem-sag}; residual regimes
and catch projections follow the same control table as the pelagics
(Figures~\ref{fig:dem-rec}--\ref{fig:dem-fcst}; \texttt{dem-f.csv}).
Resilience reading is the same as for pelagics: gaps between status-quo and
low-recruitment (or Regime) paths indicate how exposed demersal catch and
biomass are if current productivity persists or worsens further.

\begin{figure}[htbp]
  \centering
  \includegraphics[width=0.95\textwidth]{dem-sag-vs-stock.png}
  \caption{Demersal stocks: SAG versus FLStock diagnostics relative to
           reference points.}
  \label{fig:dem-sag}
\end{figure}

\begin{figure}[htbp]
  \centering
  \includegraphics[width=0.95\textwidth]{dem-rec-residuals.png}
  \caption{Demersal stock--recruit residuals with regime segments.}
  \label{fig:dem-rec}
\end{figure}

\begin{figure}[htbp]
  \centering
  \includegraphics[width=0.95\textwidth]{dem-forecast-catch.png}
  \caption{Demersal forecast catch by stock and scenario.}
  \label{fig:dem-fcst}
\end{figure}

% -----------------------------------------------------------------------------
\section{Nephrops functional units}
\label{sec:tac-neph}

Nephrops are treated as data-limited surplus-production problems. SAG biomass,
catch/landings and fishing pressure provide the observational scaffold
(Figures~\ref{fig:neph-stock}--\ref{fig:neph-f}). Of the nine functional units
carried through the resilience analysis, eight are projected here
(\texttt{nep.fu.11}, \texttt{12}, \texttt{13}, \texttt{14}, \texttt{15},
\texttt{19}, \texttt{2021}, \texttt{22}); \texttt{nep.fu.16} is excluded because
its JABBA fit did not converge, so it does not appear in \texttt{neph-f.csv}.
JABBA fits condition biodyn OMs; relative stock, yield and harvest trajectories
are shown in
Figures~\ref{fig:neph-jb-stock}--\ref{fig:neph-jb-h}. Process-error residuals
and the associated regime decomposition
(Figures~\ref{fig:neph-pe-res}--\ref{fig:neph-pe}) mirror the recruitment-regime
step used for age-structured stocks.

\begin{figure}[htbp]
  \centering
  \includegraphics[width=0.95\textwidth]{neph-stock-ts.png}
  \caption{Nephrops SAG stock biomass with $B_{\mathrm{trigger}}$ where
           available.}
  \label{fig:neph-stock}
\end{figure}

\begin{figure}[htbp]
  \centering
  \includegraphics[width=0.95\textwidth]{neph-catch-ts.png}
  \caption{Nephrops catch and landings by functional unit.}
  \label{fig:neph-catch}
\end{figure}

\begin{figure}[htbp]
  \centering
  \includegraphics[width=0.95\textwidth]{neph-f-ts.png}
  \caption{Nephrops fishing pressure relative to $F_{\mathrm{MSY}}$.}
  \label{fig:neph-f}
\end{figure}

\begin{figure}[htbp]
  \centering
  \includegraphics[width=0.95\textwidth]{neph-jabba-stock.png}
  \caption{JABBA stock biomass relative to $B_{\mathrm{MSY}}$.}
  \label{fig:neph-jb-stock}
\end{figure}

\begin{figure}[htbp]
  \centering
  \includegraphics[width=0.95\textwidth]{neph-jabba-yield.png}
  \caption{JABBA catch relative to MSY.}
  \label{fig:neph-jb-yield}
\end{figure}

\begin{figure}[htbp]
  \centering
  \includegraphics[width=0.95\textwidth]{neph-jabba-harvest.png}
  \caption{JABBA harvest relative to $F_{\mathrm{MSY}}$.}
  \label{fig:neph-jb-h}
\end{figure}

\begin{figure}[htbp]
  \centering
  \includegraphics[width=0.95\textwidth]{neph-proc-residuals.png}
  \caption{Nephrops process-error residuals with regime polygons.}
  \label{fig:neph-pe-res}
\end{figure}

\begin{figure}[htbp]
  \centering
  \includegraphics[width=0.95\textwidth]{neph-proc-error.png}
  \caption{Nephrops process-error time series.}
  \label{fig:neph-pe}
\end{figure}

Forward harvest scenarios produce catch paths relative to MSY
(Figure~\ref{fig:neph-fcst}; \texttt{neph-f.csv}). A status-quo harvest extract
is also written as \texttt{neph-tac-sq.csv}. Closed-loop MSE with a survey HCR
remains staged separately and is not part of this appendix.

\begin{figure}[htbp]
  \centering
  \includegraphics[width=0.95\textwidth]{neph-forecast-catch.png}
  \caption{Nephrops projected catch relative to MSY under the harvest-rate
           scenarios.}
  \label{fig:neph-fcst}
\end{figure}

% -----------------------------------------------------------------------------
\section{North Atlantic albacore (ICCAT)}
\label{sec:tac-alb}

Northern albacore (\texttt{alb-n}) is conditioned on the ICCAT SS3 assessment.
Assessment panels and stock--recruit residuals are shown in
Figures~\ref{fig:alb-check}--\ref{fig:alb-rec}. Advice catches bridge the stock
into the projection window; status-quo $F$ is then held forward under
recruitment multipliers of 1, 0.75 and 0.5
(Figure~\ref{fig:alb-fcst}; \texttt{iccat-f.csv}).

\begin{figure}[htbp]
  \centering
  \includegraphics[width=0.95\textwidth]{alb-assessment-panels.png}
  \caption{Albacore assessment time series with SS3 MSY-class reference
           lines (catch, SSB, $F$, recruits).}
  \label{fig:alb-check}
\end{figure}

\begin{figure}[htbp]
  \centering
  \includegraphics[width=0.75\textwidth]{alb-rec-residuals.png}
  \caption{Albacore stock--recruit residuals from the fitted Beverton--Holt
           relationship.}
  \label{fig:alb-rec}
\end{figure}

\begin{figure}[htbp]
  \centering
  \includegraphics[width=0.95\textwidth]{alb-forecast.png}
  \caption{Albacore forward projections under status-quo $F$ with recruitment
           scaled by 1, 0.75 and 0.5.}
  \label{fig:alb-fcst}
\end{figure}

% -----------------------------------------------------------------------------
\section{Reading the scenarios}
\label{sec:tac-read}

\begin{itemize}
  \item \textbf{Status-quo $F$ / harvest} isolates the effect of lower
        recruitment (or process error) while fishing mortality stays recent.
  \item \textbf{$F_{\mathrm{MSY}}$} asks whether MSY-based fishing, under the
        same recruitment assumption, recovers or erodes biomass relative to
        trigger points in the CSV ($B/B_{\mathrm{trig}}$).
  \item \textbf{Regime} (ICES age-structured stocks) substitutes the latest
        residual regime for the short-window residual mean---useful when the
        recent decade is atypical of the whole series.
  \item Absolute catch scales differ by stock; relative panels (Nephrops,
        SAG/$B_{\mathrm{trig}}$ plots) support cross-unit comparison.
\end{itemize}

\subsection{Caveats and implications}
\label{sec:tac-caveats}

\begin{itemize}
  \item \textbf{2025 SAG inputs where 2026 is incomplete.}
        Terminal states and advice bridges may shift when 2026 assessments
        are final. Implication: treat near-term levels as provisional; relative
        scenario differences are more robust than absolute tonnes.
  \item \textbf{Fishing pressure is held fixed after the bridge year.}
        There is no annual advice feedback or HCR in these runs (except the
        separate Nephrops MSE workstream). Implication: paths show exposure
        under sustained pressure, not what management would actually decide
        year by year.
  \item \textbf{Recruitment / process-error assumptions are deterministic
        scalars}, not stochastic ensembles. Implication: figures show central
        scenario contrasts, not probability intervals or risk percentages.
  \item \textbf{$B/B_{\mathrm{trigger}}$ is not on a common scale across
        Nephrops functional units.} In \texttt{neph-f.csv} it ranges from below
        $0.001$ to above $100$ between FUs (and up to $\sim 24$ for some
        demersal stocks), reflecting differences in how each FU's biomass index
        is scaled relative to its trigger. Implication: read
        $B/B_{\mathrm{trigger}}$ \emph{within} a stock (relative to $1$), not as
        a cross-FU magnitude, until the per-FU scaling is reconciled.
  \item \textbf{Figures are regenerated from the R Markdown knits.}
        Before locking V8.0, re-knit the four stock-group Rmds, refresh
        \texttt{report/latex/figs/} in R with
        \texttt{source("report/finalise/02\_stage\_figures.R"); stageFigures()},
        and confirm captions still match the panels.
\end{itemize}

% -----------------------------------------------------------------------------
\section*{Appendix data pointers}
\addcontentsline{toc}{section}{Appendix data pointers}

\begin{center}
\begin{tabular}{@{}llp{0.42\textwidth}@{}}
\toprule
Group & Forecast CSV & Notes \\
\midrule
Pelagic  & \texttt{data/TAC/csv/pel-f.csv}   & Five scenarios through 2050 (4 stocks) \\
Demersal & \texttt{data/TAC/csv/dem-f.csv}   & Same control table (8 stocks) \\
Nephrops & \texttt{data/TAC/csv/neph-f.csv}  & Harvest-rate scenarios through 2041, 8 FUs (\texttt{nep.fu.16} excluded); also \texttt{neph-tac-sq.csv} \\
Albacore & \texttt{data/TAC/csv/iccat-f.csv} & Status-quo $F$, three recruitment scales through 2050 \\
\bottomrule
\end{tabular}
\end{center}
